%---------------------------------------------------------------------------
%
%                          Vorlage der Arbeitsgruppe
%             Computer Vision and Pattern Recognition Group (CVPR)
%                           der Universität Münster
%                         http://cvpr.uni-muenster.de
%
%---------------------------------------------------------------------------
% Geeignet für:
%  - Seminararbeiten
%  - Bachelorarbeiten
%  - Masterarbeiten
%---------------------------------------------------------------------------
% Autoren:
%  - Daniel Tenbrinck
%  - Fabian Gigengack
%  - Michael Schmeing
%  - Lucas Franek
%  - Andreas Nienkötter
%---------------------------------------------------------------------------
% Version:
%  - 1.0.3 (05.10.2016)
%	 - Ersetzung von veralteten Befehlen durch Aktuelle
%	 - Einige ausführlichere Beispiele
%    - Einführung von listings
%    - Aktuelle Version der Eidesstattlichen Erklärung
%  - 1.0.2 (09.09.2011)
%    - Titelblatt um Matrikelnummer und Studiengang ergänzt
%  - 1.0.1 (05.07.2011)
%---------------------------------------------------------------------------
% 
% "THE BEER-WARE LICENSE" (Revision 42):
% The above mentioned authors wrote this file. As long as you retain this
% notice you can do whatever you want with this stuff. If we meet some day,
% and you think this stuff is worth it, you can buy us a beer in return.
%  --------------------------------------------------------------------------

\documentclass[a4paper, twoside, 12pt, ngerman]{scrbook} % Layout-Einstellungen für das Dokument

\usepackage[utf8]{inputenc} % UTF-8 Codierung
\usepackage[ngerman]{babel} % Deutsche Beschriftung

\usepackage{graphicx} % Um Bilder einzufügen
\usepackage{subfigure} % Um mehrere Bilder in eine figure einzufügen
\usepackage{amssymb, amsmath} % Für Mengensymbole und über Gleichheitszeichen schreiben
\usepackage{verbatim} % Um Quellcode in das Dokument einzufügen.
\usepackage[dvipsnames]{xcolor} % Für Farben
\usepackage[linkbordercolor=blue]{hyperref} % Für Links im Dokument
\usepackage{algorithmic} % Für Pseudo-Code
\usepackage{algorithm} % Wrapper für Pseudo-Code
\usepackage[font={small}, labelfont=bf]{caption} % kleine Bildunterschriften
\usepackage{geometry} % Für Feinanpassungen des Layouts

\usepackage{listings} % Für Code-Listings
\renewcommand{\lstlistingname}{Quelltext} %Ändert die Überschrift von Listing nach Quelltext

\usepackage{colortbl}
\usepackage{tikz}
\usetikzlibrary{tikzmark}
% Einstellungen für Abstand an den Rändern
\geometry{a4paper,left=35mm,right=35mm,top=20mm,bottom=20mm, includeheadfoot}

\begin{document}
\pagenumbering{roman}

% Titelblatt
\begin{titlepage}
\begin{centering}
\vspace*{\fill}
\includegraphics[width=12cm]{./img/wwu-logo-neu.pdf}

\vspace{3cm} 

{\LARGE
	\textbf{Titel der Arbeit\\
			(mehrzeilig auch möglich)}\\[1.2cm]
}

{\large
	Bachelorarbeit oder Masterarbeit oder Seminararbeit\\[2cm]
}

{\large
	vorgelegt von:
}

{ \Large
	\textbf{Vorname Nachname}\\[1cm]
}

{\large
	Matrikelnummer: 123456\\[2mm]
}

{\large
	Studiengang: [Studiengang]\\[1cm]
}
    
{\large
	Thema gestellt von:
}

{\Large
	\textbf{Prof. Dr. Xiaoyi Jiang}\\[1cm]
}
                               
{\large
	Arbeit betreut durch:
}

{\Large
	\textbf{Vorname Nachname}\\[1cm]
}

{\large
Münster, \today
}
\vfill
\end{centering}
\end{titlepage}

% Abstract und Danksagung
\chapter*{Zusammenfassung}

An dieser Stelle soll eine Zusammenfassung der wissenschaftlichen Arbeit stehen. Idealerweise ist die Zusammenfassung genau eine Seite lang, was in etwa 350 - 450 Wörtern entspricht. Sie bildet zusammen mit dem Fazit einen Rahmen der gesamten wissenschaftlichen Arbeit und sollte auch für Außenstehende einfach zu verstehen und ansprechend formuliert werden.

% Inhaltsverzeichnis
\tableofcontents

\cleardoublepage
\pagenumbering{arabic}

% Die Hauptkapitel der Arbeit
\chapter{Einleitung}
\label{ch:einleitung}

Im Einleitungskapitel werden dem Leser das Thema, die Grundlagen und vergleichbare Arbeiten präsentiert.

Im Folgenden werden ein paar wichtige Tipps zum Erstellen der wissenschaftlichen Arbeit festgehalten. Die vorliegende inhaltliche Strukturierung dient nur zur Orientierung und ist nicht verbindlich. Insbesondere bei Seminararbeiten kann die Struktur abweichen. Beispielsweise ist dort nicht generell eine Implementierung gefordert (und folglich müssen auch keine Ergebnisse beschrieben werden). Auch Zusammenfassung, Danksagung, eidesstattliche Erklärung, Abbildungs- und Tabellenverzeichnis können bei Seminararbeiten weggelassen werden.


\section{Referenzierung von Unterkapiteln}
\label{s:ref} 
Man sollte direkt jedes Kapitel, Unterkapitel und jede Formel mit einem Label versehen \verb+\label{}+ um eine konsistente Referenzierung im gesamten Dokument zu ermöglichen. Eine Referenzierung im Fließtext lässt sich mittels \verb+\ref{}+ umsetzen.


\section{Zeilenumbrüche und Absätze}
\label{s:zeilenumbruch} 
Zeilenumbrüche sollten im Quelltext immer mittels einer Leerzeile umgesetzt werden, um eine automatische Texteinrückung in der darauffolgenden Zeile zu ermöglichen. Dies macht das Lesen der Arbeit einfacher. Auf die Verwendung des LaTeX-Kommandos \verb+\\+ sollte verzichtet werden.


\section{Einfügen von Grafiken}
\label{s:grafik} 
Grafiken sollten mittels der \verb+\figure+-Umgebung eingebettet werden. Um die Druckqualität und Wiederverwendbarkeit der Grafiken für Vorträge, Poster, etc.\ zu erhöhen sind Vektorgrafiken (z.B.\ .eps oder .pdf) zu bevorzugen. Zur Erzeugung und Konvertierung von Vektorgrafiken ist die OpenSource Software \emph{Inkscape} zu empfehlen.

Um mehrere Bilder horizontal anzuordnen, sollte die \verb+\subfigure+-Umgebung verwendet werden. Diese Bilder können dann mit \verb+\subref+ oder \verb+\ref+ referenziert werden und erscheinen im Text als \subref{fig:subfig1} oder \ref{fig:subfig1}.

\begin{figure}
\begin{center}
	\subfigure[Das Logo]{
		\includegraphics[width=0.45\columnwidth]{./img/wwu-logo-neu.pdf}
		\label{fig:subfig1}
	}
	\subfigure[Noch einmal das Logo]{
		\includegraphics[width=0.45\columnwidth]{./img/wwu-logo-neu.pdf}
		\label{fig:subfig2}
	}
\end{center}
\caption{Das Logo der WWU}
\label{fig:histogram}
\end{figure}


\section{Wissenschaftliches Zitieren}
\label{s:zitat}
Für das Referenzieren von wissenschaftlicher Literatur wie Fachbüchern, Konferenzpapern und Veröffentlichungen in Wissenschaftsmagazinen gibt es unterschiedliche Konventionen. Je nach Fachrichtung weichen Layout und Zitationsstil sehr stark voneinander ab \cite{jele2010}. Wir empfehlen aus Gründen der Einheitlichkeit die Verwendung der \emph{Bibtex}-Umgebung. Die zitierte Literatur kann ausgelagert in einer Datei (z.B.: Quellen.bib) gepflegt werden und mittels des \verb+\cite+-Kommandos referenziert werden.


\section{Zusammenfassung}
\label{s:zusammenfassung}
Zum Ende eines längeren Kapitels bietet es sich häufig an eine Zusammenfassung der wichtigsten Punkte zu liefern. Dies erleichtert das Lesen und den Übergang zum nächsten Kapitel.
\chapter{Methodik}
\label{ch:methodik}

Im Methodik-Kapitel werden die mathematischen Ausführungen des Verfahrens bzw. des Algorithmus vorgestellt, jedoch technische Details zur Umsetzung und Implementierung (falls nötig) auf ein darauffolgendes Kapitel verlagert.


\section{Mathematische Notation}
\label{s:notation}

Mathematische Formeln können mittels der \verb+\begin{align}...\end{align}+ Umgebung gesetzt werden:

\begin{align}
f(n) & =
	\begin{cases}
		n/2, & \text{wenn }n\text{ gerade,}\\
		3n+1, & \text{wenn }n\text{ ungerade.}
	\end{cases}
\label{eq:f} \\
%
g(n) & = \frac{n}{2} \label{eq:g}
\end{align}

\section{Algorithmus}
\label{s:algorithmus}

Eigene Algorithmen beschreibt man am Besten mit Hilfe von Pseudo-Code und dem Paket \verb+algorithm+.

\begin{algorithm}
\caption{Algorithmus}
\label{alg:alg}
\begin{algorithmic}
\algsetup{indent=2em}

\REQUIRE Argument $n\in\mathbb{N}$
\STATE $a = 0$
\FOR{ $i=0,\dots,n$}
	\STATE $a = a + 1$
\ENDFOR
\RETURN $a$
\end{algorithmic}
\end{algorithm}
\chapter{Implementierung}
\label{ch:implementierung}

In diesem Kapitel werden Technische Details, wie zum Beispiel die Verwendete Laufzeitumgebung, Laufzeitanalyse oder wichtige Implementierungsdetails behandelt.
Dabei sollte beachtet werden inwiefern diese Angaben für die Arbeit von Bedeutung sind. Die Namen einzelner Funktionsaufrufe oder Klassendiagramme sind zum Beispiel im Allgemeinen für den Leser uninteressant.

\section{Quelltext-Abschnitte}
Eigene Auszüge aus dem Quelltext bindet man am Einfachsten mit dem Befehl \verb+\lstinputlisting{}+ ein. Das Ergebnis ist in \ref{fig:code} zu sehen.
Optionen für die Sprache, Tab-Breite etc. von \verb+\lstinputlisting+ können auch am Anfang des Dokuments mit \verb+\lstset{}+ gesetzt werden.
Das Paket erlaubt mit \verb+firstline=...+ und \verb+lastline=...+ auch die Einbindung von einzelnen Zeilen einer Datei.

%Auch möglich: Kapselung in \begin{figure} ... \end{figure} für flexible Positionierung
\lstinputlisting[language=C,tabsize=2,frame=single,caption={Codebeispiel}]{src/myTest.m}
\label{fig:code}



\chapter{Ergebnisse}
\label{ch:ergebnisse}

Dieses Kapitel sollte die Ergebnisse beinhalten, die mit den Methoden aus \autoref{ch:methodik} erstellt wurden.

\begin{table}[h]
\caption{Beispieltabelle}
\begin{center}
	\begin{tabular}{|c||c|c|}
		\hline
		Spalte1 & Spalte2 & Spalte3 \\ 
		\hline\hline
		   1    &    2    &    3    \\ 
		\hline
	\end{tabular}
\end{center}
\label{tbl:table}
\end{table}
\chapter{Diskussion}
\label{ch:diskussion}

In diesem Kapitel werden die zuvor vorgestellten Ergebnisse der Arbeit diskutiert. Häufig wird es mit dem Fazit zusammengelegt.
\chapter{Fazit}
\label{ch:fazit}

Dieses Kapitel bildet die abschließende Zusammenfassung der Arbeit. Dazu können die folgende Punkte behandelt werden:
\begin{itemize}
	\item Reflexion: wurden die Ziele der Arbeit erreicht?
	\item mögliche Erweiterungen und Verbesserungen (\glqq future work\grqq)
\end{itemize}

% Abbildungs- und Tabellenverzeichnis
\listoffigures
\listoftables

% Literaturverzeichnis
\bibliographystyle{unsrtdin}
\bibliography{Quellen}

% Eidesstattliche Erklärung
\chapter*{Eidesstattliche Erklärung}

% Die Aktuelle Version der Eidesstättlichen Erklärung kann beim zuständigen Prüfungsamts gefunden werden.
% Nachfolgend ist die Version für Bachelor und Master des Prüfungsamts Mathe/Informatik vom 05.10.2016

Hiermit versichere ich, dass die vorliegende Arbeit über \textit{\glqq Titel\grqq} selbstständig verfasst worden ist, dass keine anderen Quellen und Hilfsmittel als die angegebenen benutzt worden sind und dass die Stellen der Arbeit, die anderen Werken – auch elektronischen Medien – dem Wortlaut oder Sinn nach entnommen wurden, auf jeden Fall unter Angabe der Quelle als Entlehnung kenntlich gemacht worden sind.

\vspace{1cm}

\parbox{20em}{\hrulefill}

Vorname Nachname, Münster, \today

\vspace{1cm}

Ich erkläre mich mit einem Abgleich der Arbeit mit anderen Texten zwecks Auffindung von Übereinstimmungen sowie mit einer zu diesem Zweck vorzunehmenden Speicherung der Arbeit in eine Datenbank einverstanden.

\vspace{1cm}

\parbox{20em}{\hrulefill}

Vorname Nachname, Münster, \today

\end{document}
